% Akhil Indurti -- CV.  Build: pdflatex cv.tex   (see README.md)
\documentclass[10pt,letterpaper]{article}
\usepackage[left=0.5in,textwidth=7.4in,top=0.33in,bottom=0.5in]{geometry}
\usepackage{tabularx}
\usepackage[colorlinks,allcolors=blue,pdftitle={Akhil Indurti},pdfauthor={Akhil Indurti}]{hyperref}

% ---- Type and page ------------------------------------------------------
\makeatletter
\renewcommand\normalsize{\@setfontsize\normalsize\@viiipt{10pt}}   % 8pt Computer Modern on a 10pt grid
\makeatother
\normalsize
\pagestyle{empty}
\setlength\parindent{0pt}
\setlength\parskip{0pt}
\raggedbottom
\hyphenation{ScyllaDB MongoDB DigitalOcean}   % product names never hyphenate

% ---- Layout lengths ------------------------------------------------------
\newlength\bodyindent  \setlength\bodyindent{0.75in}  % body column is inset this much from both ends of the rules
\newlength\tablehang   \setlength\tablehang{0.45in}   % label tables (Skills, Projects) stick out this far left of the body
\newlength\labelgap    \setlength\labelgap{12pt}      % gap between a table label and its text (= 2\tabcolsep)
\newlength\rowsep      \setlength\rowsep{5pt}         % extra space between table rows
\newlength\sectionskip \setlength\sectionskip{8pt}    % extra space above a section heading
\newlength\rulegap     \setlength\rulegap{2.2pt}      % heading baseline to rule
\newlength\afterrule   \setlength\afterrule{12pt}     % rule to first body line
\newlength\entrysep    \setlength\entrysep{6pt}       % extra space between entries
\newlength\bulletsep   \setlength\bulletsep{6pt}      % extra space above a bullet list

% ---- Macros --------------------------------------------------------------
\makeatletter
\newcommand\sectionrule{\par\nobreak\vspace{\rulegap}\hrule\nobreak}

% \cvheader{Name}{contact lines separated by \\} -- the contact block is right-aligned
% under the rule, and the first section heading shares the baseline of its third line.
\newcommand\cvheader[2]{%
  \textbf{#1}\sectionrule
  \vspace{\baselineskip}%
  \hfill\smash{\begin{tabular}[t]{@{}r@{}}#2\end{tabular}}\par
  \addvspace{\baselineskip}}

% \begin{cvsection}{HEADING} ... \end{cvsection}: flush-left bold heading with a full-width
% rule; everything inside is set in the indented body column.
\newenvironment{cvsection}[1]{%
  \par\addvspace{\sectionskip}%
  \textbf{#1}\sectionrule
  \vspace{\afterrule}\nobreak
  \begin{list}{}{%
    \setlength\leftmargin{\bodyindent}%
    \setlength\rightmargin{\bodyindent}%
    \setlength\itemsep{\entrysep}%
    \setlength\topsep{0pt}\setlength\partopsep{0pt}\setlength\parsep{0pt}%
    \setlength\labelwidth{0pt}\setlength\labelsep{0pt}\setlength\itemindent{0pt}%
    \interlinepenalty\@M}% page breaks only between entries
}{\end{list}}

% \entry{Title}{Location}{Organisation}{Dates}: a bold title line with the organisation
% (slanted) and dates flush right. Leave {Location} or {Organisation} empty to drop them.
\newcommand\entry[4]{%
  \item\textbf{#1}\if\relax\detokenize{#2}\relax\else{} @ #2\fi
  \hfill\mbox{\if\relax\detokenize{#3}\relax\else\textsl{#3}, \fi#4}\par\nobreak}

% Bullet list under an entry.
\newenvironment{bullets}{%
  \begin{list}{--}{%
    \setlength\leftmargin{22pt}\setlength\labelsep{5pt}%
    \setlength\labelwidth{\dimexpr\leftmargin-\labelsep\relax}%
    \renewcommand\makelabel[1]{\hfil##1}%
    \setlength\topsep{\bulletsep}\setlength\partopsep{0pt}%
    \setlength\itemsep{1pt}\setlength\parsep{0pt}%
    \@beginparpenalty\@M \@itempenalty\@M}}
  {\end{list}}

% \begin{labeled} \row{Label}{text} ... \end{labeled}: right-aligned bold labels, wrapped text.
% Each row ends with \rowsep of space; the environment's end cancels the last one.
\newenvironment{labeled}{%
  \item[]\hspace*{-\tablehang}%
  \tabularx{\dimexpr\linewidth+\tablehang\relax}[t]{@{}>{\bfseries}r@{\hspace{\labelgap}}X@{}}}
  {\noalign{\vskip-\rowsep}\endtabularx\par}
\newcommand\row[2]{#1 & #2\tabularnewline\noalign{\vskip\rowsep}}

\makeatother
\newcommand\gh[1]{\href{https://github.com/smasher164/#1}{#1}}   % one of my GitHub repos

% ==========================================================================
\begin{document}

\cvheader{Akhil Indurti}{%
  \textbf{\href{mailto:aindurti@gmail.com}{aindurti@gmail.com} (703-431-5466)}\\
  \href{https://www.akhil.cc}{akhil.cc}\\
  \href{https://github.com/smasher164}{github.com/smasher164}}

\begin{cvsection}{SKILLS}
\begin{labeled}
\row{Languages}{C, Go, Rust, C++, Assembly, Java, Python, JS, SML, OCaml, Raku, C\#, F\#, Racket, Shell, SQL, \LaTeX.}
\row{Tools}{HLS, SCTE-35, Widevine DRM, LLVM \& Go toolchains, OpenGL, Git, Docker, Kubernetes, Linux, Nix, GraphQL, Android,
  Adobe Creative Suite, QEMU, Azure, AWS, GCP, DigitalOcean, fly.io, Terraform, ScyllaDB,
  MongoDB, Concourse, ZFS.}
\end{labeled}
\end{cvsection}

\begin{cvsection}{PROJECTS}
\begin{labeled}
\row{\href{https://go-review.googlesource.com/q/owner:aindurti}{golang/go}}{Member who contributes
  patches, triages issues, and reviews code. Works on the Go compiler, standard library, and runtime.
  Notable contributions: fused-multiply-add, static analysis checks, and playground.\par
  \textbf{Speaker at GopherCon 2021}: Gave a 45 min talk about
  \href{https://www.gophercon.com/agenda/session/593026}{``Writing a Static Analyzer for Go Code''}.}
\row{\gh{mem}}{A memory allocator for Go.}
\row{\gh{nas}}{Configuration for Home NAS, running NixOS, ZFS, Tailscale, Home Assistant, Samba, and Plex.}
\row{\gh{oscons}}{Notes on OS Construction.}
\row{\gh{mexdown}}{An extensible human-readable plain-text format.}
\row{\gh{arewesixelyet}}{Site lists support for SIXEL graphics across terminals. Got to front-page of HN.}
\row{\gh{pw-volume}}{Interface to PipeWire volume controls, before WirePlumber had support. Over 60 active users.}
\row{\gh{gflat}}{An ML-like programming language with rank-n polymorphism, typeclasses, algebraic data types,
  green threads, and data-race freedom.}
\row{\gh{burrito}}{Build system and package manager on an ML-like language with lexical effect handlers.}
\row{\gh{tapl}}{Typed lambda calculus solutions for learning type theory from \textit{Types and Programming Languages}.\par
  \textbf{Participant at OPLSS 2021}: Participated in the
  \href{https://www.cs.uoregon.edu/research/summerschool/summer21/participants.php}{Oregon Programming Languages Summer School}
  to study interesting topics in PL Theory and Formal Verification.}
\end{labeled}
\end{cvsection}

\begin{cvsection}{EXPERIENCE}

\entry{Engineer III, Software Process Engineering}{Los Angeles, CA}{Samsung Ads}{May, 2026 - Present}
Maintain the Origin Manifest Manipulator, Pipeline, and Sessions services that deliver Samsung TV Plus content and handle
server-side ad insertion (SSAI).
\begin{bullets}
\item Migrated the legacy Origin service into Pipeline, and rolled out Transmit as an alternative SSAI
  provider to AWS Elemental MediaTailor.
\item Descaled the Sessions service to cut infrastructure costs, and built dashboards and telemetry for
  stream and ad-insertion health.
\end{bullets}

\entry{Software Engineer II on Azure}{Redmond, WA}{Microsoft}{March, 2023 - Nov, 2025}
Service Fabric team: System Services (Failover Manager, Transport, Federation) and Infrastructure Services.
\begin{bullets}
\item Worked on routing, failover, networking, and VM upgrades for the distributed systems platform that
  underpins much of Azure.
\item On call for live-site incidents, driving mitigations and RCAs.
\end{bullets}

\entry{Software Engineer II on .NET}{Redmond, WA}{Microsoft}{May, 2022 - March, 2023}
\begin{bullets}
\item \textbf{APICompat}: Worked on the public rollout for a tool that detects binary breaking changes between two
  .NET assemblies.
\item \textbf{GenAPI}: Generates C\# source from an assembly. The generated source contains stub
  declarations for types and members, where declarations preserve the API contract of the
  original source. Fixed Roslyn bugs for assembly metadata.
\item \textbf{Doc source porting tool}: Spearheaded internal effort to port externally maintained XML
  docs into inline docs for the .NET repos.
\end{bullets}

\entry{Software Engineer II}{Mount Laurel, NJ}{Comcast}{May, 2020 - April, 2022}
Working on Comcast's Next Generation Access Network (NGAN) to build Genome, a poller for millions of
network devices across the country.
\begin{bullets}
\item Implemented an SNMP connection pool that significantly sped up the long tail of our polling cycle,
  and cut AWS costs.
\item Led a testing \& deployment strategy to discover test devices and monitor health statuses.
\item Detected incompatible changes to GraphQL and Avro schemas.
\item Automated datalake topic generation and IAM roles via Terraform.
\item Assisted in MongoDB to ScyllaDB migration, and multi-region setup.
\item Improved and hired talent to improve code quality and CI infrastructure.
\end{bullets}

\entry{Research Assistant}{Charlottesville, VA}{ShiftLab}{Jan, 2020 - May, 2020}
Worked under Professor Samira Khan to find crash-consistency bugs via structured fuzzing, as well add support
for persistent memory to the Go programming language. Contributed to:
\begin{bullets}
\item \href{https://dl.acm.org/doi/abs/10.1145/3373376.3378452}{Cross-Failure Bug Detection in Persistent Memory Programs}
\item \href{https://dl.acm.org/doi/abs/10.1145/3445814.3446691}{PMFuzz: test case generation for persistent memory programs}
\end{bullets}

\entry{Teaching Assistant}{Charlottesville, VA}{University of Virginia}{Spring 2018 - Spring 2019}
\begin{bullets}
\item \textbf{Theory of Computation}: Course covers automata theory, complexity theory, and problem-solving
  techniques. I assist in grading, holding office hours, and supervising problem-solving sessions. I developed
  a new \href{https://www.tmsim.akhil.cc}{Turing Machine Simulator assignment}.
\item \textbf{Program Data and Representation}: Course covers C++, data structures, and assembly. I assist in
  grading, holding office hours, and running lab. I developed a new
  \href{https://uva-cs.github.io/pdr/labs/lab05/index.html}{Tree lab} that introduces AVL trees.
\end{bullets}

\entry{Research}{Charlottesville, VA}{Center for Automata Processing}{Fall 2017 - Spring 2018}
I worked to explore the application and implementation of state machines. Under my advisors Professor
Kevin Skadron and Jack Wadden, \textit{my research} focused on developing and implementing an NFA reduction
algorithm to perform prefix and suffix merging in VASim.

\entry{Android SDK}{Palo Alto, CA}{TRNQL}{January - August 2015}
Worked to enable easy integration of contextual information into an application. Created example
applications for clients to study, and led meetups to introduce the community to these SDK features.

\entry{Software Engineering Intern}{Milpitas, CA}{FireEye}{Summer 2014}
Implemented a web application for the management and education on the company's security products.
\begin{bullets}
\item Wrote a web application to educate, display, and manage the new intrusion prevention system.
\item Developed a real-world understanding of network security in the context of large organizations.
\end{bullets}

\end{cvsection}

\begin{cvsection}{EDUCATION}

\entry{University of Virginia (B.S. in Computer Science)}{}{}{}
\textbf{Coursework}: Operating Systems, Advanced Computer Architecture,
Graduate Compilers, Machine Learning, Computer Graphics, Capstone Research,
Algorithms, Theory of Computation, Program Data and Representation

\end{cvsection}

\end{document}
